\section*{Erbsensuppe mit Würstchen}
\ihead{}\chead{Personen:4}\ohead{}
\ifoot{Vorbereitungszeit:10 min}\cfoot{}\ofoot{Kochzeit:30 min}
{\Large Zutaten}
\begin{multicols}{2}
\begin{itemize}
\item \num{2} große Kartoffeln, mehligkochend
\item \num{1} Paket TK-Erbsen, ca. \SI{300-400}{g}
\item \num{4} Wiener Würstchen
\item \num{1} Zwiebel
\item \SI{1}{l} Gemüsebrühe
\item etwas Salz und Pfeffer
\item \SI{1-2}{EL} Öl
\item etwas Kondensmilch, 10,% Fett
\end{itemize}
% \columnbreak
\end{multicols}
\noindent {\Large Zubereitung}\ \
Bei diesem Rezept müssen die Erbsen vorab nicht aufgetaut werden.
Die Zwiebel schälen und in Würfel schneiden.
Die Kartoffeln schälen und ebenfalls in kleine Würfel schneiden.
In einem Topf ca. \SI{1-2}{EL} Öl erhitzen und die Zwiebeln darin schön braun anbraten.
Die Kartoffeln dazugeben und mit etwa der Hälfte der Gemüsebrühe ablöschen.
Einmal aufkochen lassen.
Jetzt die Erbsen zugeben und so viel Brühe hinzugeben, dass das Gemüse schön bedeckt ist und schwimmen kann.
Nun die Suppe ca. \SI{15}{min} köcheln lassen.
Das Gemüse ist gar, wenn sich die Kartoffeln problemlos zerdrücken lassen.
Den Topf vom Herd nehmen.
Etwas Kondensmilch zugeben und alles mit einem Pürierstab cremig pürieren.
Wer mag, kann auch vorher ein oder zwei Schöpflöffel Gemüse beiseite geben, dann hat man danach noch Stückchen in der Suppe.
Die Würstchen in Scheiben schneiden und zur Suppe geben.
Jetzt den Topf wieder auf den Herd geben und nochmals erwärmen.
Mit Salz und Pfeffer abschmecken.
